\section{Архитектурный коллапс}
\label{sec:collapse}

Этот раздел излагает и доказывает два результата, на которых держится весь
проект. Оба суть \status{Т}: они опираются на установленные теоремы УГМ вместе
с одним элементарным диофантовым фактом и теоремой Гурвица. Ничто здесь не
постулируется.

\subsection{Единственная инцидентная структура}

Пусть $\mathsf{N}$ обозначает инцидентность плоскости Фано $\Fano$: семь точек
$\{A,S,D,L,E,O,U\}$, семь прямых
\[
\{A,S,L\},\ \{A,D,E\},\ \{A,O,U\},\ \{S,D,O\},\ \{S,E,U\},\ \{D,L,U\},\ \{L,E,O\},
\]
каждая прямая несёт три точки, каждая точка лежит на трёх прямых, и всякая пара
точек лежит ровно на одной общей прямой (система Штейнера $S(2,3,7)$). Пишем
$k^\ast(i,j)$ для \emph{медиатора} пары --- единственной третьей точки их общей
прямой.

\begin{theorembox}[Теорема~\ref{thm:collapse} (Архитектурный коллапс)]
На семимодовом УГМ-субстрате четыре архитектурных слоя --- интерконнект, код
коррекции ошибок, алгебра инструкций и правило обучения --- суть один объект:
каждый есть функция единственной инцидентности $\mathsf{N}$.
\end{theorembox}

\begin{theorem}\label{thm:collapse}
Пусть состояние холона есть $\Gam\in\Dens(\mathbb{C}^7)$, эволюционирующее под
УГМ-линдбладианом $\LindOmega$. Тогда:
\begin{enumerate}[leftmargin=1.7em,itemsep=1pt]
  \item[\textbf{(L1)}] \emph{Интерконнект.} Физическая связь --- семь прямых
        $\mathsf{N}$: сектор $i$ достигает сектора $j$ только через медиатора
        $k^\ast(i,j)$ (закон замыкания: когерентность $\gamma_{ij}$ создаётся или
        чинится только живыми плечами $\gamma_{i k^\ast},\gamma_{k^\ast j}$; две
        когерентности рождают третью тогда и только тогда, когда тройка --- прямая).
  \item[\textbf{(L2)}] \emph{Код.} Проверочная матрица кода Хэмминга $[7,4,3]$
        имеет столбцами семь ненулевых векторов $\mathbb{F}_2^3$; при любой
        разметке это семь точек $\mathsf{N}$, а шестнадцать кодовых слов включают
        ровно семь слов веса 3, чьи носители --- семь прямых. Инцидентность
        интерконнекта \emph{есть} проверочная геометрия кода.
  \item[\textbf{(L3)}] \emph{Алгебра инструкций.} Кубическая (нелинейная) часть
        $\LindOmega$ порождается октонионными структурными константами $f_{ijk}$
        --- единственной $G_2$-инвариантной трилинейной формой на $\mathrm{Im}(\Oct)$
        --- ненулевыми ровно когда $\{i,j,k\}$ --- прямая (правило отбора Фано).
        Таблица умножения инструкции есть линейная инцидентность $\mathsf{N}$.
  \item[\textbf{(L4)}] \emph{Обучение.} Регенеративный член $\Regen$ гонит $\Gam$
        к $\rhostar=\vp(\Gam)$, и весь поток ковариантен под
        $G_2=\mathrm{Aut}(\Oct)$, чьё действие на семь мод имеет $\mathsf{N}$
        орбитной инцидентностью ($G_2$-ригидность). Правило неподвижной точки
        (обучения) есть функция той же $\mathsf{N}$.
\end{enumerate}
Следовательно $\mathrm{L1}=\mathrm{L2}=\mathrm{L3}=\mathrm{L4}$ как функции
$\mathsf{N}$. \hfill$\square$
\end{theorem}

\begin{corollary}[Нулевой налог рассогласования]\label{cor:zerotax}
Поскольку четыре слоя суть одна структура, нет межслойных интерфейсов для
согласования. Площадь, статическая мощность и усилие верификации, которые
обычная машина тратит на примирение сеть~$\leftrightarrow$~код,
код~$\leftrightarrow$~ISA и ISA~$\leftrightarrow$~оптимизатор
(Таблица~\ref{tab:four-layers}), в \TALOS{} тождественно ноль: есть одна
структура, которую надо разместить, верифицировать и тактировать.
\end{corollary}

\begin{figure}[H]
\centering
\begin{tikzpicture}[font=\small,>=Stealth]
  \node[draw,rounded corners,fill=talossteel!8,minimum width=2.7cm,minimum height=0.7cm] (net)  at (0,3)   {Интерконнект (NoC)};
  \node[draw,rounded corners,fill=talossteel!8,minimum width=2.7cm,minimum height=0.7cm] (ecc)  at (0,2)   {Код ошибок (ECC)};
  \node[draw,rounded corners,fill=talossteel!8,minimum width=2.7cm,minimum height=0.7cm] (isa)  at (0,1)   {Инструкция (ISA)};
  \node[draw,rounded corners,fill=talossteel!8,minimum width=2.7cm,minimum height=0.7cm] (lrn)  at (0,0)   {Правило обучения};
  \node[align=center,text=plosgray] at (0,3.9) {\textbf{Сегодня: четыре подсистемы}\\\scriptsize проектируются порознь, налог на каждом шве};
  \begin{scope}[shift={(6.4,1.5)}]
    \coordinate (fA) at (0,1.7); \coordinate (fS) at (-1.5,-0.9); \coordinate (fD) at (1.5,-0.9);
    \coordinate (fL) at (-0.75,0.4); \coordinate (fE) at (0.75,0.4); \coordinate (fO) at (0,-0.9); \coordinate (fU) at (0,0);
    \draw[talosamber,thick] (fA)--(fS) (fA)--(fD) (fS)--(fD) (fA)--(fO) (fS)--(fE) (fD)--(fL);
    \draw[talosamber,thick] (0,-0.034) circle[radius=0.866];
    \foreach \n in {fA,fS,fD,fL,fE,fO,fU}{ \fill[talossteel] (\n) circle (2.4pt); }
    \node[font=\scriptsize\bfseries,above] at (fA) {A};
    \node[font=\scriptsize\bfseries,left]  at (fS) {S};
    \node[font=\scriptsize\bfseries,right] at (fD) {D};
    \node[font=\scriptsize\bfseries,fill=white,inner sep=0.5pt] at (-0.95,0.5) {L};
    \node[font=\scriptsize\bfseries,fill=white,inner sep=0.5pt] at (0.95,0.5)  {E};
    \node[font=\scriptsize\bfseries,below] at (fO) {O};
    \node[font=\scriptsize\bfseries,fill=white,inner sep=0.5pt] at (0.22,0.18) {U};
  \end{scope}
  \node[align=center,text=talosamber] at (6.4,3.9) {\textbf{\TALOS{}: один объект}\\\scriptsize плоскость Фано $\mathsf{N}$};
  \foreach \b in {net,ecc,isa,lrn}{ \draw[->,plosgray,thick] (\b.east) to[out=0,in=180] (4.4,1.5); }
  \node[text=plosgray,font=\scriptsize] at (3.4,0.4) {коллапс};
\end{tikzpicture}
\caption{Архитектурный монизм. В каждой прежней машине четыре слоя --- отдельные
подсистемы, соединённые на облагаемых налогом швах (слева).
Теорема~\ref{thm:collapse} доказывает, что на семимодовом субстрате они суть
одна инцидентная структура --- плоскость Фано (справа) --- поэтому интерфейсы, и
их налог, исчезают (Следствие~\ref{cor:zerotax}).}
\label{fig:collapse}
\end{figure}

\begin{zerobox}[С нуля --- один рисунок, который есть четыре игры]
Начерти мелом на асфальте классики. Один и тот же рисунок --- четыре разные
вещи разом, смотря что ты с ним делаешь. Он --- \emph{поле}: куда можно
ставить ногу (сеть). Он --- \emph{правило}: какие прыжки законны
(инструкция). Он --- \emph{судья}: неверный прыжок виден мгновенно, нога вне
линии (проверка). И он --- \emph{учитель}: играя по сетке снова и снова,
становишься лучше (обучение). Никто не свинчивал четыре прибора; один
меловой рисунок делает все четыре работы \emph{в силу своей формы}.
Теорема~\ref{thm:collapse} говорит: для семичастной само-наблюдающей машины
существует ровно один такой рисунок --- семилинейная фигура справа на
Рис.~\ref{fig:collapse}, --- а Теорема~\ref{thm:unique} --- что ни у какого
другого числа частей его нет.
\end{zerobox}

\subsection{Почему ровно семь --- и нигде более}

Скептик согласится, что \emph{если} $N=7$, то четыре слоя совпадают, и спросит,
особенна ли семёрка или лишь одна удобная точка. Она уникальна, и причина
резче, чем сознание-теоретические доказательства минимальности УГМ: это
арифметический тай-брейк, решаемый принципом третьего порядка.

\begin{theorembox}[Теорема~\ref{thm:unique} (Уникальность коллапса)]
Архитектурный коллапс Теоремы~\ref{thm:collapse} реализуем при единственном
числе узлов: $N=7$.
\end{theorembox}

\begin{theorem}\label{thm:unique}
Потребуем, чтобы три целочисленных условия выполнялись при одном $N$:
\begin{itemize}[leftmargin=1.6em,itemsep=0pt]
  \item[\textbf{(C1)}] $N$ --- порядок подлинной проективной плоскости:
        $N=q^2+q+1$ для степени простого $q\ge 2$ \emph{(так что интерконнект ---
        симметричный дизайн $S(2,q{+}1,N)$ из $N$ шин-прямых)};
  \item[\textbf{(C2)}] $N$ --- совершенная длина Хэмминга: $N=2^r-1$
        \emph{(так что инцидентность --- совершенный однократно-исправляющий код)};
  \item[\textbf{(C3)}] $N$ --- размерность мнимой части нормированной алгебры с
        делением с \emph{ненулевым ассоциатором} \emph{(так что алгебра
        инструкций несёт подлинную третьепорядковую, неассоциативную структуру)}.
\end{itemize}
Тогда $N=7$ единственно.
\end{theorem}

\begin{proof}
\textbf{(C3) одна вынуждает $7$.} По теореме Гурвица нормированные алгебры с
делением суть $\mathbb{R},\mathbb{C},\mathbb{H},\Oct$ размерностей $1,2,4,8$, с
мнимыми частями размерности $0,1,3,7$. Из них лишь $\Oct$ неассоциативна (её
ассоциатор $[x,y,z]\ne 0$); остальные ассоциативны, поэтому их трилинейная форма
вырождена и третьепорядковой структуры нет. Значит (C3) $\Rightarrow N=7$.

\textbf{(C1)$\cap$(C2) повторяется, но лишь $7$ переживает (C3).} Диофантово
уравнение $q^2+q+1=2^r-1$ имеет целочисленные решения $q\in\{2,5,90\}$, т.е.\
$N\in\{7,31,8191\}$ (исчерпывающе для $q\le 200$). Из них лишь $q=2,5$ суть
степени простого, поэтому лишь $N=7,31$ суть порядки \emph{подлинных}
проективных плоскостей ($q=90$ не степень простого, поэтому $\mathrm{PG}(2,90)$
не галуасова плоскость, а $N=8191$ --- голое целочисленное совпадение, не
интерконнект). Итак совпадение сеть-\emph{есть}-код повторяется при $N=7$ и
$N=31$ --- семёрка не единственное место, где фабрика может быть кодом Хэмминга.
Но $\{7,31,8191\}\cap\{0,1,3,7\}=\{7\}$: из всякого совпадения лишь семь есть
размерность мнимой части алгебры с делением, и (по абзацу выше) лишь семь
\emph{неассоциативна}. Пересечение всех трёх условий есть $\{7\}$. \end{proof}

\begin{remark}[Тай-брейк --- это принцип третьего порядка]
Содержание Теоремы~\ref{thm:unique} в том, что ассоциатор --- дискриминатор.
Среди чисел узлов, где фабрика может \emph{быть} кодом ($7,31,8191$), лишь семь
позволяет инструкции подлинно жить на \emph{тройках}, а не на парах --- принцип
третьего порядка УГМ, доказывающий, что парная структура слепа к Фано и что
структура начинается с трёх. При $N=31$ или $8191$ можно всё ещё слить
интерконнект и ECC, но алгебра инструкций коллапсирует к ассоциативной (парной),
а правило обучения теряет кубический генератор: четырёхкратный коллапс
проваливается. Вырожденный кандидат $N=3$ (кватернионы) исключён дважды ---
кватернионы ассоциативны, а $q=1$ не проективная плоскость. Семёрка --- единственное
целое, при котором сеть, код, неассоциативная инструкция и ковариантная
неподвижная точка суть одна вещь.
\end{remark}

\begin{remark}[Связь с доказательствами минимальности УГМ]
УГМ уже доказывает $N=7$ из сознания (теорема минимальности семимерности) и из
алгебры (октонионный вывод и $G_2$-ригидность). Теорема~\ref{thm:unique} ---
\emph{третий, независимый} путь к тому же целому, родной для компьютерной
архитектуры: она просит лишь, чтобы сеть машины была её кодом, а инструкция ---
третьепорядковой, и выводит семь. То, что три несвязанных требования ---
жизнеспособный разум, алгебра с делением и самокорректирующийся третьепорядковый
компьютер --- падают на одно число, само есть сильнейшее свидетельство того, что
семь структурно, а не выбрано. \status{Т}
\end{remark}
